\documentclass{article}

% Submission/preprint build. For a workshop, swap to e.g.
%   \usepackage[dblblindworkshop]{neurips_2025} \workshoptitle{...}
% and add `final` only after acceptance.
\usepackage[preprint]{neurips_2025}

\usepackage[utf8]{inputenc}
\usepackage[T1]{fontenc}
\usepackage{hyperref}
\usepackage{url}
\usepackage{booktabs}
\usepackage{amsfonts}
\usepackage{amsmath}
\usepackage{nicefrac}
\usepackage{microtype}
\usepackage{xcolor}
\usepackage{multirow}
\usepackage{graphicx}

\newcommand{\typ}[1]{\texttt{#1}}

\title{Beyond the Ledger: Cross-Source Multi-Agent Fraud Auditing\\of ERP Journal Entries and Invoice Documents}

% TODO: replace with real author names and affiliations.
\author{%
  Author Name$^{1}$ \\
  $^{1}$Department, Institution, City, Country \\
  \texttt{author@institution.edu}
}

\begin{document}

\maketitle

\begin{abstract}
Recent work shows that large language models (LLMs) can act as anomaly detectors over journal entries, but such detectors read the ledger alone. Many procurement frauds are not visible in a single ledger row: an altered invoice amount, a duplicated invoice number or a vendor missing from the master file only appears when the ledger is checked against documents and against other entries. We present a multi-agent auditing framework in which transaction, document and reconciliation agents each produce evidence, a stacked meta-classifier fuses their scores, and a local LLM is confined to narrating that evidence under a numeric grounding check. We evaluate on the 533,009-row synthetic SAP journal-entry ledger from the deepAD repository of \citet{schreyer2017detection}, linked to 20,000 rendered invoice images and injected with five fraud typologies, under a temporal, vendor-disjoint split. The fused system reaches a test PR-AUC of 0.854 (95\% CI 0.840--0.869) and an F1 of 0.911 at a recall of 0.997; removing the reconciliation agent lowers PR-AUC to 0.727. A fully label-free variant of the agents reaches 0.826, indicating that most of the gain comes from cross-source evidence rather than from supervision. Re-implemented on the same rows with a 3B-parameter model, the ledger-only AuditCopilot prompt of \citet{kadir2025auditcopilot} reaches a PR-AUC of only 0.614 against a 0.601 base rate, because it follows an Isolation Forest hint that is wrong on most genuine rows; adding our agents' evidence sharpens its discrimination but leaves it far below the same evidence scored by rule. We report the negative results alongside the positive ones: nearly all remaining false positives trace to a labelling artifact in the benchmark, several rules mirror the injected typologies, and our SHAP attributions fail a deletion-faithfulness test.
\end{abstract}

\section{Introduction}

Auditors test journal entries to find errors and fraud, and the standard practice, rule-based Journal Entry Tests (JETs) \citep{droste2024jet}, produces large numbers of false positives that must be reviewed by hand \citep{kadir2025auditcopilot}. Unsupervised learning has long been applied to the problem: deep autoencoders \citep{schreyer2017detection}, adversarial autoencoders \citep{schreyer2019adversarial}, and clustering with Isolation Forest \citep{herreros2025applied,gronewald2024anomaly}. Most recently, AuditCopilot \citep{kadir2025auditcopilot} showed that a prompted LLM, given dataset statistics and an Isolation Forest hint, can outperform JETs and classical baselines while explaining its decisions in natural language.

All of these detectors read the ledger and nothing else. In procurement, however, several common fraud schemes are defined by a disagreement \emph{between} sources. An altered invoice carries a different amount from its ledger posting; a duplicate invoice reuses an invoice number already posted; a ghost vendor is paid but is absent from the vendor master. A detector that sees a single ledger row in isolation has, by construction, no access to the field that differs. Our question is therefore not whether an LLM can read a journal entry, but which evidence an auditing system needs, and how much of a detector's performance depends on having it.

We make four contributions:
\begin{itemize}
  \item A multi-agent framework with transaction, document and reconciliation agents, fused by an out-of-fold stacked meta-classifier \citep{wolpert1992stacked}. The LLM only narrates evidence the agents already computed, and a grounding check rejects any narrative containing a number that is not in that evidence.
  \item A linked benchmark: the 533,009-row synthetic SAP ledger from the deepAD repository of \citet{schreyer2017detection}, extended with vendors, dates, invoice and purchase-order numbers, line descriptions and 20,000 rendered invoice images, with five injected typologies and a temporal, vendor-disjoint split.
  \item A same-data comparison against AuditCopilot and Isolation Forest, including a label-free variant of our agents that separates the value of cross-source evidence from the value of supervision.
  \item An explicit account of what does not work: the remaining false positives come from a benchmark labelling artifact, a generator artifact initially made our document agent useless, and our SHAP explanations fail a deletion-faithfulness test.
\end{itemize}

\section{Method}

Each agent returns an evidence envelope: a risk score in $[0,1]$ and a list of named findings with the values that triggered them. Downstream components, including the reviewer dashboard and the LLM, only ever see these envelopes.

\paragraph{Transaction agent.} A gradient-boosted classifier \citep{chen2016xgboost} over twelve engineered numeric features and six categorical SAP fields (currency, company code, transaction key, profit centre, posting key, GL account). The numeric features include the amount's z-score relative to its vendor and to its account, its ratio to the vendor median, vendor posting velocity over seven days, days since the vendor's previous posting, posting-to-document lag, a round-amount flag and the vendor's first-digit (Benford) deviation. Vendor and account baselines are fit on genuine training rows only. An Isolation Forest \citep{liu2008isolation} trained on genuine rows supplies its anomaly score as an additional feature, and the output is calibrated with isotonic regression on the validation split.

\paragraph{Document agent.} Signals computable from one invoice in isolation: the invoice amount's z-score within its line-item description class (graded linearly between 2 and 6 standard deviations, fit on genuine training rows), a multi-scale round-amount check, a missing purchase-order number, a vague description, and low extraction confidence. The score is a weighted sum capped at 1. Checks that need the ledger or a second document belong to the reconciliation agent, so the same evidence is not counted twice in fusion.

\paragraph{Reconciliation agent.} Eight deterministic rules, each logged by name when it fires: invoice amount differs from the ledger amount by more than 3\%; vendor name fuzzy-match score below 85; the invoice number appears on another posting anywhere in the full ledger; the vendor is absent from the vendor master; the purchase-order number is missing; posting and document dates are more than 20 days apart; a low-value line description is posted at an amount above the 90th percentile of genuine training amounts; and a \emph{split purchase}, where two or more postings from the same vendor on the same purchase order fall within seven days, each below the approval limit, while their total exceeds it. The score is the sum of fixed rule weights, capped at 1. The duplicate and split-purchase checks are cross-posting tests, so they scan the entire ledger, not just the document-linked subset, and cite the other postings' document numbers, dates and amounts. The ledger carries no approval-policy table, so the limit is estimated as the 98th percentile of all training amounts, which requires no labels.

\paragraph{Fusion.} A class-balanced logistic regression takes the three agent scores as input. To avoid the most common stacking error, the transaction scores used to train it are out-of-fold: the transaction agent is retrained five times, each time with one fold of the meta-training rows held out. Decision thresholds (F1-optimal, and the highest threshold reaching 90\% recall) are chosen on the validation split only. The learned coefficients are 14.06 for reconciliation, 3.97 for transaction and 2.32 for document.

\paragraph{Coordinator.} A local LLM (Qwen2.5-3B-Instruct \citep{qwen2025qwen25}, temperature 0, fixed seed) receives the agent scores, the fired rules and the extracted document fields, but never the label, and writes a short finding with a recommended action for the reviewer (prompt in Appendix~\ref{app:prompts}). The prompt separates the entry under review from any other entry a rule cites, because without that separation the model attached a duplicated posting's document number and date to the entry under review. A grounding check extracts every number from the narrative and rejects it, with one retry, if any number does not appear in the supplied evidence. The fraud decision is always made by the fused score; the LLM never decides.

\paragraph{Review interface.} Cases are served on a dashboard showing the invoice image, agent scores, the full ledger row, fired rules and the narrative; for duplicate invoices it shows a field-by-field comparison with the duplicated posting and both invoice images. Reviewer decisions are written to an append-only audit log.

\section{Dataset and Experiments}

\paragraph{Ledger.} We use the full synthetic SAP journal-entry dataset (533,009 rows) from the public deepAD repository (\url{https://github.com/GitiHubi/deepAD}) accompanying the autoencoder work of \citet{schreyer2017detection}. We add synthetic vendors (306), posting and document dates spanning January 2025 to June 2026, invoice and purchase-order numbers, and line-item descriptions whose value tier is derived from the posting amount, with 22\% of tiers redrawn at random so that the description is informative about the amount without being a deterministic function of it.

\paragraph{Fraud typologies.} We inject five typologies for a combined fraud rate of 2.06\% (10,998 rows): \typ{threshold\_split} (2,470; a high-value posting and 2--3 same-vendor, same-week neighbours are rewritten so that each amount falls just under an approval threshold, set at the 98th amount percentile, while their sum exceeds it), \typ{duplicate\_invoice} (2,132; a posting reuses another posting's invoice number and vendor), \typ{ghost\_vendor} (2,132; a posting assigned to a vendor absent from the master file), \typ{desc\_amount\_mismatch} (2,132; a posting above the 90th amount percentile relabelled with a low-value description), and \typ{altered\_amount} (2,132; the invoice amount differs from the ledger amount by 15--60\%).

\paragraph{Split and documents.} Vendors are ordered by median posting date and assigned to train, validation and test as whole blocks (310,972 / 92,717 / 129,320 rows), so the split is temporal and no vendor appears in more than one partition. We render 20,000 invoice images with layout templates and scan-like degradation: every fraudulent row, the genuine original of every duplicated invoice, and a random genuine sample. Imaging every row would require about 110\,GB. Fusion is evaluated on the 4,852 document-linked test rows, of which 60.1\% are fraudulent; \textbf{a random scorer therefore has a PR-AUC of 0.601 on this subset}, and all PR-AUC values below should be read against that base rate.

\paragraph{Experiments.} On the full document-linked test subset we report the transaction-only and document-only agents, the full framework, single-agent ablations (refitting the meta-classifier without one agent) and per-typology recall, each with 1,000-sample bootstrap 95\% confidence intervals.

For the comparison with AuditCopilot, an LLM call per row is too slow on our hardware (about 19\,s per call on a CPU-only laptop) to cover the full subset, so we draw a stratified sample of 400 test rows (40 per typology and 200 genuine), weight every pooled metric back to the subset's composition, and bootstrap within strata. All methods are scored on exactly the same rows:
\begin{itemize}
  \item \textbf{Isolation Forest}, AuditCopilot's ML baseline: 200 trees, default contamination, fit on all training rows with vendor and account baselines computed without labels.
  \item \textbf{AuditCopilot}: the prompt of \citet[Fig.~1]{kadir2025auditcopilot} with dataset statistics, the Isolation Forest verdict and score, and the entry's ledger fields. We made three changes, each forced by our data: the ledger has no posting-user field, so ``user'' statistics refer to vendors; the model is Qwen2.5-3B-Instruct, the largest that runs at a usable speed on our hardware; and vendor identifiers are pseudonymised, because the injected ghost-vendor identifiers carry a \texttt{GHOST-} prefix that would give the label away in plain text.
  \item \textbf{AuditCopilot + signals}: the same prompt with our document and reconciliation agents' fired evidence appended. Comparing it with plain AuditCopilot measures whether cross-source evidence helps an LLM detector, independent of our fusion model.
  \item \textbf{Agents (label-free)}: document and reconciliation scores combined by noisy-OR, $1-(1-s_d)(1-s_r)$, with every baseline fit on all training rows without labels, flagged at 0.30 (the weight of one strong rule).
  \item \textbf{Fusion (ours)}: the full supervised framework, at its validation-selected F1 threshold.
\end{itemize}

\section{Results and Discussion}

\begin{table}[t]
  \caption{Same-data comparison on a stratified 400-row sample of the document-linked test subset, weighted to its composition (base-rate PR-AUC 0.601). R: recall per typology at each method's own decision rule; FPR: false-positive rate on genuine rows. LLM methods (Qwen2.5-3B-Instruct) emit a binary verdict, so their PR-AUC reflects a single operating point.}
  \label{tab:comparison}
  \centering
  \small
  \setlength{\tabcolsep}{3.5pt}
  \begin{tabular}{lccccccccccc}
    \toprule
    Method & Labels & PR-AUC [95\% CI] & P & R & F1 & R\textsubscript{dup} & R\textsubscript{ghost} & R\textsubscript{split} & R\textsubscript{desc} & R\textsubscript{alter} & FPR \\
    \midrule
    Isolation Forest      & no  & 0.698 [0.653, 0.747] & 0.62 & 0.85 & 0.72 & 0.75 & 0.98 & 0.90 & 0.95 & 0.65 & 0.78 \\
    AuditCopilot          & no  & 0.614 [0.598, 0.629] & 0.61 & 0.91 & 0.73 & 0.83 & 0.90 & 1.00 & 1.00 & 0.83 & 0.86 \\
    AuditCopilot + signals & no & 0.630 [0.618, 0.642] & 0.63 & 0.99 & 0.77 & 1.00 & 0.98 & 1.00 & 1.00 & 1.00 & 0.88 \\
    Agents (label-free)   & no  & 0.826 [0.790, 0.865] & 0.81 & 0.87 & 0.84 & 1.00 & 1.00 & 1.00 & 0.25 & 1.00 & 0.31 \\
    Fusion (ours)         & yes & \textbf{0.853} [0.814, 0.891] & \textbf{0.83} & \textbf{1.00} & \textbf{0.91} & 1.00 & 1.00 & 1.00 & 1.00 & 1.00 & 0.31 \\
    \bottomrule
  \end{tabular}
\end{table}

\begin{table}[t]
  \caption{Agent-level results and ablations on the full document-linked test subset (4,852 rows, base-rate PR-AUC 0.601), with 1,000-sample bootstrap 95\% CIs. Ablations refit the meta-classifier without the named agent. At its validation-selected threshold the full system reaches precision 0.839, recall 0.997 and F1 0.911 (2,907 TP, 556 FP, 10 FN).}
  \label{tab:ablation}
  \centering
  \small
  \begin{tabular}{llcc}
    \toprule
    & Configuration & PR-AUC [95\% CI] & ROC-AUC [95\% CI] \\
    \midrule
    \multirow{2}{*}{Single agent} & Transaction only & 0.688 [0.670, 0.708] & 0.649 [0.633, 0.665] \\
                                  & Document only    & 0.663 [0.646, 0.680] & 0.594 [0.580, 0.608] \\
    \midrule
    \multirow{3}{*}{Ablation} & w/o document       & 0.859 [0.844, 0.872] & 0.859 [0.848, 0.871] \\
                              & w/o transaction    & 0.837 [0.822, 0.851] & 0.844 [0.831, 0.856] \\
                              & w/o reconciliation & 0.727 [0.710, 0.747] & 0.656 [0.640, 0.672] \\
    \midrule
    Full & Fusion (ours) & 0.854 [0.840, 0.869] & 0.855 [0.843, 0.868] \\
    \bottomrule
  \end{tabular}
\end{table}

\paragraph{Main comparison.} On the shared sample (Table~\ref{tab:comparison}), the fused system reaches a PR-AUC of 0.853 and an F1 of 0.91, against 0.698 and 0.72 for Isolation Forest, which flags 78\% of genuine rows. Isolation Forest's high recall is therefore not useful: under the vendor-disjoint split every test vendor is unseen in training, so vendor-relative features shift and most test rows look anomalous.

\paragraph{AuditCopilot on our data.} The re-implemented AuditCopilot reaches a PR-AUC of 0.614 [0.598, 0.629], barely above the 0.601 base rate (ROC-AUC 0.525), because it flags 86\% of genuine rows. Its high recall, including 0.83 on \typ{altered\_amount}, whose evidence is on the invoice it never sees, is a consequence of flagging nearly everything rather than of detection. The Isolation Forest hint drives this. When the hint reads \emph{anomaly}, the LLM flags 97\% of genuine and 96\% of fraudulent rows alike, and under our vendor-disjoint split the hint reads \emph{anomaly} on 78\% of genuine rows. \citet{kadir2025auditcopilot} found the hint essential for precision; on our data the hint is wrong most of the time, and the LLM follows it.

Cross-source evidence does help the LLM where the hint leaves room. On rows whose hint reads \emph{normal}, the gap between the flag rates of fraudulent and genuine rows grows from 14 points (61\% vs.\ 48\%) for plain AuditCopilot to 35 points (97\% vs.\ 61\%) when our agents' evidence is added, and recall on \typ{duplicate\_invoice} and \typ{altered\_amount} rises from 0.83 to 1.00. Overall, however, the \emph{+ signals} variant reaches only 0.630, far below the same evidence scored by rule without labels (0.826). The evidence matters more than its reader, and a 3B-parameter model is a poor reader: it does not follow the prompt's instruction to flag only clear anomalies. A larger model may do so; we could not run one at a usable speed on our hardware (Section~\ref{sec:limitations}).

\paragraph{Is the gain from supervision?} A supervised system beating zero-shot detectors would be unremarkable. The label-free agents reach a PR-AUC of 0.826 [0.790, 0.865], whose interval overlaps that of the supervised fusion (0.853 [0.814, 0.891]). Most of the advantage over ledger-only baselines therefore comes from cross-source evidence, not from labels. Supervision helps mainly on \typ{desc\_amount\_mismatch} (recall 0.25 to 1.00), whose rule weight (0.22) falls below the label-free threshold; the meta-classifier learns to weight it.

\paragraph{Which evidence matters.} Removing the reconciliation agent causes the largest loss (PR-AUC 0.854 to 0.727, ROC-AUC 0.855 to 0.656; Table~\ref{tab:ablation}). Removing the document agent does not lower PR-AUC (0.859, overlapping CI): its signal is largely redundant with reconciliation, which flags the same rows by rule. We consider ``independently informative'' the supportable claim for the document agent, not ``improves the ensemble''.

\paragraph{Recall near 1.00 is not the result.} All five typologies reach a recall of 0.98 or higher under fusion, but this should not be read as detection skill. The rules that catch them encode the same definitions the injectors used: a duplicated invoice number, a vendor absent from the master file, a ledger--invoice amount difference far beyond the 3\% tolerance, postings that share a purchase order, and, for \typ{desc\_amount\_mismatch}, a rule that is close to the injector's inverse. What these numbers establish is that the evidence is \emph{necessary}: for \typ{altered\_amount} the differing field exists only on the invoice, so no ledger-only detector can see it, whatever its model. They do not establish that the rules generalise to real fraud.

\paragraph{Split purchases need a cross-posting test.} Without the split-purchase rule, \typ{threshold\_split} was detected at a recall of only 0.204 and the full system's PR-AUC was 0.789: each piece of a split purchase is unremarkable on its own, so no single-row signal can see it. The rule raises that recall to 0.984 on the full test subset while firing on 2 of 126,403 genuine test postings. It is, however, sensitive to the estimated approval limit: estimated from genuine training rows only (5.04M rather than 5.19M), the limit falls below the largest split pieces, and recall drops to 0.338. In practice the limit is a known policy value, not an estimate. Isolation Forest's 0.90 recall on this typology comes with a 78\% false-positive rate and is not a usable alternative.

\paragraph{Where the remaining errors are.} Of the 556 false positives on the test subset, 544 are the genuine originals of duplicated invoices, which the duplicate rule necessarily flags along with their copies. The benchmark cannot separate them: the injector chooses which posting of a pair is the fraudulent one without regard to time, so the fraudulent copy is the later posting in only 50.3\% of the 2,115 pairs. In practice a duplicate is submitted after its original, and a rule could flag only the later posting. Removing the duplicate rule's weight from the genuine originals, as a projection rather than a run on a regenerated corpus, would raise precision from 0.839 to 0.995. We report this as a benchmark limitation rather than a result. For a reviewer, the pair is the finding, and the dashboard already presents both postings side by side.

\paragraph{A generator artifact, not a model failure.} Our document agent initially scored a test ROC-AUC of 0.500 and received a \emph{negative} fusion coefficient ($-3.48$). Rewriting the agent did not help. The cause was the data generator, which assigned each description's value tier from a hash of the GL transaction type, so descriptions carried no information about amounts (0.05 standard deviations between class means; Spearman correlation $-0.003$). After tiers were derived from amounts, the lift of the amount-versus-description feature rose from 1.17 to 5.47, the agent's ROC-AUC to 0.594, and its fusion coefficient to $+2.96$. We report this because synthetic fraud benchmarks can silently remove the signal an agent is designed to find.

\paragraph{Explainability.} SHAP attributions \citep{lundberg2017unified} for the transaction agent are stable across retraining seeds (mean top-5 Jaccard overlap 0.962), and the dominant features are vendor frequency rank, the vendor-relative amount z-score and the Isolation Forest score. They are not, however, faithful by a deletion test: replacing each row's top-3 SHAP features with their medians lowers the calibrated score by 0.000 on average, against 0.0014 for three random features. We therefore do not claim faithful feature attributions. The auditor-facing explanation in our system is the list of fired rules and their values, which is faithful by construction, and the LLM narrative is constrained to restate it.

\section{Conclusion}

Ledger-only detectors, whether classical or LLM-based, cannot see fraud whose evidence lives in another source. A multi-agent design that checks the ledger against invoice documents and against other postings recovers most of that signal, even without labels, while the LLM remains a constrained narrator rather than the decision-maker. On shared data (Table~\ref{tab:comparison}), a small-model re-implementation of AuditCopilot barely exceeds the base rate, and supplying it with our agents' evidence helps it discriminate but does not close the gap: the evidence is worth more read by rule than read by a small LLM. Two of the most effective signals, duplicate invoices and split purchases, are tests across postings rather than properties of one row. Future work should regenerate the benchmark with time-ordered duplicates, replace ground-truth document fields with real OCR, test typology variants that the rules were not written for, and evaluate on labelled real-world ledgers at their natural fraud rate.

\section*{Limitations}
\label{sec:limitations}

\paragraph{Synthetic data and rule alignment.} All fraud is injected, and the reconciliation rules mirror the injectors' definitions, including the shared purchase order that the split-purchase rule keys on, so per-typology recall overstates what the same rules would achieve on real fraud. The split-purchase rule was added after an error analysis on the test split; its sensitivity to the estimated approval limit is reported above. Our split is vendor-disjoint and temporal, which is stricter than a random split, but external validity on real ledgers remains to be shown.

\paragraph{No OCR in the evaluated pipeline.} The document and reconciliation agents currently read the invoice fields from the corpus's ground truth rather than from OCR. Two attempts to run PaddleOCR \citep{du2020ppocr} over the rendered images failed, one on runtime and one on a CUDA library conflict. Results therefore reflect perfect extraction; the extraction-confidence signal never fires, and the effect of OCR noise is unmeasured.

\paragraph{Enriched evaluation subset.} Evaluation uses the document-linked subset, which is 60.1\% fraud and contains every duplicated invoice's genuine original. Operating-point precision does not transfer to the corpus's natural 2.06\% fraud rate, where, with a 31\% false-positive rate on genuine rows, it would be substantially lower; we did not measure it there.

\paragraph{LLM baseline fidelity.} We re-implement AuditCopilot with a 3B-parameter model on CPU, smaller than the 7--8B models that gave the original paper's best results (its 2B model also produced many false positives). Our model over-flags (86\% of genuine rows), so our numbers likely understate what AuditCopilot achieves with a stronger model, and they should be read as a comparison with a small-model re-implementation, not with the original system. The 400-row sample also gives wide intervals. LLM output at temperature 0 was not reproducible across model reloads on our hardware, so every response is cached and the cache is the artifact our numbers are computed from.

\paragraph{Deployment and governance.} The system is decision support. The LLM narrates but never decides, every reviewer action is logged, and any deployment would require human oversight, access controls and validation on real audit data.

\bibliographystyle{plainnat}
\bibliography{references}

\appendix

\section{Related Work}

\paragraph{Machine learning for journal-entry anomalies.} \citet{schreyer2017detection} detect anomalous journal entries with deep autoencoders, and their public deepAD repository provides the synthetic SAP ledger we build on; \citet{schreyer2019adversarial} extend this with adversarial autoencoders whose latent space supports interpretation. Clustering and Isolation Forest have been applied to purchase processes \citep{herreros2025applied} and to synthetic general-ledger data \citep{gronewald2024anomaly}. Tree ensembles remain strong on tabular data \citep{grinsztajn2022tree}, which motivates our transaction agent.

\paragraph{LLMs for anomaly detection and auditing.} AnoLLM \citep{tsai2025anollm} and AD-LLM \citep{yang2024adllm} study LLMs as tabular anomaly detectors. AuditCopilot \citep{kadir2025auditcopilot} applies prompted LLMs to tax-related ledgers, combining dataset statistics with an Isolation Forest hint; it is the base paper we compare against. AuditAgent \citep{bai2025auditagent} uses expert-guided multi-agent reasoning to discover fraud evidence across financial filings. Our work differs from both in operating on ERP postings linked to invoice documents, and in using agents to produce verifiable evidence rather than LLM verdicts.

\paragraph{Explainability in auditing.} \citet{gnoss2022xai} explain an autoencoder anomaly detector for audit through a surrogate model evaluated with auditors. We use SHAP \citep{lundberg2017unified} and report a deletion-faithfulness test, which our attributions fail.

\section{Prompt Templates}
\label{app:prompts}

\paragraph{Coordinator.} The narrative prompt used by our coordinator; braces denote fields filled from the agent envelopes.
\begin{quote}\small\ttfamily
You are a fraud audit coordinator. Summarize the evidence below into a finding of at most \{max\_words\} words for a human reviewer. State only what the evidence supports; do not speculate, invent numbers, or guess at causes not listed here.\\[2pt]
Case: row \{row\_id\}, typology under review: \{typology\}\\[2pt]
The entry under review. Every field here describes THIS entry: \{doc\_fields\}\\[2pt]
Agent scores (0=low risk, 1=high risk): transaction\_agent, document\_agent, reconciliation\_agent, fused risk score (tier)\\[2pt]
Reconciliation rules fired. These compare the entry above against the ledger and against other entries. Where a rule names a second posting, that posting's document number, date and amount belong to the OTHER entry -- never restate them as if they were the entry under review: \{recon\_evidence\}\\[2pt]
Write the finding, then a one-line recommended action.
\end{quote}

\paragraph{AuditCopilot re-implementation.} We use the template of \citet[Fig.~1]{kadir2025auditcopilot} with ``user'' replaced by ``vendor'' throughout. The \emph{+ signals} variant appends, before the final instruction:
\begin{quote}\small\ttfamily
CROSS-SOURCE EVIDENCE (invoice document and ledger reconciliation checks):\\
- Invoice document: vendor '\{name\}', amount \{amount\}, invoice \{invoice\_number\}\\
- \{rule\}: \{detail\} \quad (one line per fired rule, or ``No document or reconciliation check fired.'')
\end{quote}

\section{Reproducibility}

All numbers are produced by scripts in the accompanying repository with a single global seed (42): \texttt{run\_fusion.py} (fusion), \texttt{run\_experiments.py} (Table~\ref{tab:ablation}), \texttt{run\_xai.py} (SHAP stability and faithfulness), and \texttt{run\_auditcopilot\_comparison.py} (Table~\ref{tab:comparison}), whose LLM responses are cached keyed by model and prompt hash.

\end{document}
